% sec-03: the medical oncology evidence, led by the drug.
\section{Pharmacologic Rationale and the Prior Record}

Daraxonrasib (RMC-6236) is a RAS(ON) multi-selective inhibitor. In May 2026,
RASolute 302 reported median overall survival of \textbf{13.2 months} against
\textbf{6.6 months} for chemotherapy in the RAS G12 population of previously
treated metastatic pancreatic cancer \cite{rasolute302}. Nine months earlier, a
ten-arm quantitative systems pharmacology simulation with 250 ordinary
differential equations per patient had returned \textbf{12.8} against
\textbf{5.4 months} for a daraxonrasib combination \cite{qspmetpancre}.

The proximity of those two results is a chronology observation and a
hypothesis-supporting one. It is \textbf{not} a validation claim. Three
differences are material and are stated wherever the comparison appears: sample
size, 1000 simulated against 241 enrolled; intervention, a daraxonrasib
combination against daraxonrasib as a single agent; and molecular selection,
KRAS G12C against a primarily G12D and G12V population.

\begin{apptable}
\begin{tabularx}{\textwidth}{>{\raggedright\arraybackslash}p{3.4cm} >{\raggedright\arraybackslash}p{1.9cm} >{\raggedright\arraybackslash}p{1.9cm} >{\raggedright\arraybackslash}X}
\headrow \thead{Source and date} & \thead{Test arm} & \thead{Control} & \thead{Reading}\\
\midrule
QSP simulation, Aug 2025 & 12.8 mo mOS & 5.4 mo mOS & Best arm HR 0.25; 6 of 7 archetypes below HR 1.00 \\
Digital twin, Oct 2025 & 12.1 mo mOS & not applicable & Best mPFS HR 0.31; credibility 81.9 under ASME V\&V 40 \\
RASolute 302, May 2026 & 13.2 mo mOS & 6.6 mo mOS & Independent readout, RAS G12 metastatic population \\
\end{tabularx}
\end{apptable}
\tabcap{The simulated and observed results side by side. The 2025 simulated
2.4-fold survival ratio and the 2026 observed 2.0-fold ratio are reported as a
chronology and not as a validation of the model.}

None of that evidence speaks to the resectable setting, which is precisely why
the consortium question is worth asking. The metastatic readout tells the field
that the pathway is druggable in humans; it does not say what a
curative-intent resection specimen looks like after perioperative exposure.
